
\section{Pielou's Third Diversity Axiom} \label{sec:pielou-third-axiom}

The last of Pielou's three axioms relates to the behavior of a diversity measure when multiple, independent demographic attributes are considered simultaneously.

\textbf{Additivity under independent attributes.} 
``\emph{In a multiple classification (two or more qualitative variates) where the classifications are independent, the total diversity should equal the sum of the respective diversities.}'' {\cite{pielou1977mathematical}}.  
This statement can be formalized as follows, interpreting ``independent'' in the sense of probability distributions: suppose that, for a given population $\pop = \standardTriple \in \Pops_{nd\ell}$, there exist two other populations $\pop_1 = (\mX, \mU \in \R^{n\times \ell_1}, \mu)$ and $\pop_2 = (\mX, \mV \in \R^{n\times \ell_2}, \mu)$ such that $\ell = \ell_1\ell_2$ and a bijection $\alpha: [\ell_1] \times [\ell_2] \rightarrow [\ell]$ such that $y_{ij} = u_{i\alpha_1(j)} v_{i\alpha_2(j)}$ for all $i \in [n]$ and $j \in [\ell]$. 
Then $f(\pop) = f(\pop_1) + f(\pop_2)$.

This axiom, in combination with Pielou's other two and under the mild assumption of separability and continuity, is sufficient to uniquely characterize the Shannon entropy: 
\begin{thm}
    If a function $\psi: \simplex{\ell} \rightarrow \R$ is strictly concave, continuous\footnote{Convex functions are continuous on the interior of their domain, so this stipulation requires only that they are also continuous at the boundary.}, separable,  and satisfies the additivity under independent attributes property, then\footnote{Throughout this article, we will use the notation $\log$ to denote the natural logarithm, as well as the standard convention that $0 \log 0 = 0$.}  
    \begin{align}
      \psi(\vy) = -c\sum_{j = 1}^{\ell} y_j \log y_j \; \label{eq:shannon-entropy}
    \end{align}
    for some constant $c > 0$.
\end{thm}
A proof is given is given by Chakrabarti and Chakrabarty {\cite{chakrabarti2005shannon}}, where additivity under independent attributes corresponds to the authors' Axiom 4.\footnote{The authors' statement of their result also hypothesizes nonnegativity of $\phi$, but this hypothesis is not required by their proof.}







\section{Supplementary Proofs} \label{sec:supplementary-proofs}

\begin{proof}[Proof of \Cref{thm:monotone-under-group-addition}]
    Let $\phi$ be strictly convex. 
    We'll first assume that $\phi$ is nonnegative. 
    Then, for any $I \subset [\ell]$ and $\lambda = \frac{|I|}{\ell}$, we have
    \begin{align*}
      \phi\paren{\frac{1}{\ell} \ve_{[\ell]} } 
      & = \phi\paren{\lambda \frac{1}{|I|} \ve_{I} + (1-\lambda) \frac{1}{\ell - |I|} \ve_{[\ell]\setminus I}} \\
      & < \lambda \phi\paren{\frac{1}{|I|} \ve_{I}} + (1-\lambda) \phi\paren{\frac{1}{\ell - |I|} \ve_{[\ell]\setminus I}} \\
      & \leq \phi\paren{\frac{1}{|I|} \ve_{I}} \;,
    \end{align*}
    where the first inequality is by convexity and the second is by nonnegativity. 

    Now, relax the requirement that $\phi$ is nonnegative. 
    Since the simplex $\simplex{\ell}$ is compact (closed and bounded) and $\phi$ is continuous, $\phi$ attains a minimum on $\simplex{\ell}$, which we denote $\phi_{\min} = \min_{\vy \in \simplex{\ell}} \phi(\vy)$.
    Then, let $\phi'(\vy) = \phi(\vy) - \phi_{\min}$, which is nonnegative and convex. 
    The calculation above now goes through replacing $\phi$ with $\phi'$, proving that $\phi'\paren{\frac{1}{\ell} \ve_{[\ell]}} \leq \phi'\paren{\frac{1}{|I|} \ve_{I}}$  and therefore that $\phi\paren{\frac{1}{\ell} \ve_{[\ell]}} \leq \phi\paren{\frac{1}{|I|} \ve_{I}}$.
  \end{proof}

\begin{proof}[Proof of \Cref{thm:diversity-axioms-1}]
      To demonstrate maximization under the uniform, for any vector $\vy \in \simplex{\ell}$,
      \begin{align*}
        \ve_{[\ell]}/\ell = \frac{1}{\abs{\cS}}\sum_{\sigma \in \cS} \sigma(\vy) \;,
      \end{align*}
      where $\cS$ is the set of all permutations of $[\ell]$.
      It follows that 
      \begin{align*}
        \psi\paren{\frac{1}{\ell} \ve_{[\ell]}} 
        & = \psi\paren{\frac{1}{\abs{\cS}}\sum_{\sigma \in \cS} \sigma(\vy)} \\
        & < \frac{1}{\abs{\cS}}\sum_{\sigma \in \cS} \psi\paren{\sigma(\vy)} = \psi(\vy) \;,
      \end{align*}
      where the inequality is by an inductive application of strict concavity and the final equality is by symmetry.
  \end{proof}


\begin{proof}[Proof of \Cref{thm:exchanges}]
  The proof is by direct calculation. Let $\pop' = \exc{a}{b}{i}{j}{\epsilon}(\pop)$, and let $\vy_i'$ and $\vy_j'$ be the demographic vectors at locations $i$ and $j$ in $\pop'$. 
  Since the global demographic vector $\bar{\vy}$ is unchanged by the exchange, we have
  \begin{align*}
    I_{\phi}(\pop) - I_{\phi}(\pop') &= \sum_{i = 1}^n \brackets{\mu_i \phi(\vy_i) - \mu_i' \phi(\vy_i')} \\ 
    &= \mu_i\brackets{ \phi(\vy_i) - \phi(\vy_i')} + \mu_j\brackets{ \phi(\vy_j)  - \phi(\vy_j')} \\ 
    &= \mu_i \nabla \phi(\vy_i)^{\top} (\vy_i - \vy_i') \\ 
    & + \mu_j \nabla \phi(\vy_j)^{\top} (\vy_j - \vy_j') \\ 
    &+ o(\norm{\vy_i - \vy_i'}) + o(\norm{\vy_j - \vy_j'})\;. 
  \end{align*}
  since $\vy_k = \vy_k'$ for all $k \notin \{i, j\}$ and $\mu_k = \mu_k'$ for all $k$.
  Using now that $\vy_i - \vy_i' = \lambda_i (\tilde{\ve}_a - \tilde{\ve}_b)$ and $\vy_j - \vy_j' = \lambda_j (\tilde{\ve}_b - \tilde{\ve}_a)$ and ignoring the higher-order terms, we have 
  \begin{align*}
    I_{\phi}(\pop) - I_{\phi}(\pop') &\simeq \mu_i \lambda_i \nabla \phi(\vy_i)^{\top} (\tilde{\ve}_a - \tilde{\ve}_b) + \mu_j \lambda_j \nabla \phi(\vy_j)^{\top} (\tilde{\ve}_b - \tilde{\ve}_a)\;. 
  \end{align*}
  By the separability assumption, we have $\nabla \phi(\vy) = (\phi'(y_1), \ldots, \phi'(y_{\ell}))^{\top}$, and so 
  \begin{align*}
    I_{\phi}(\pop) - I_{\phi}(\pop') &\simeq \mu_i \lambda_i (\phi'(y_{ia}) - \phi'(y_{ib})) + \mu_j \lambda_j (\phi'(y_{jb}) - \phi'(y_{ja})) \\ 
    &= \epsilon (\phi'(y_{ia}) - \phi'(y_{ib}) + \phi'(y_{jb}) - \phi'(y_{ja})) \\ 
    &= \epsilon (\phi'(y_{ia}) - \phi'(y_{ja}) + \phi'(y_{jb}) - \phi'(y_{ib})) \\ 
    &> 0\;.
  \end{align*}
  The last line follows from the assumptions that $y_{ia} < y_{ja}$ and $y_{ib} > y_{jb}$, and the fact that, if $\phi$ is strictly convex, then $\phi'$ is strictly increasing, and so $\phi'(y_{ia}) - \phi'(y_{ja}) > 0$ and $\phi'(y_{jb}) - \phi'(y_{ib}) > 0$.
  Choosing $\epsilon$ sufficiently small to justify the linear approximation completes the proof. 
\end{proof}

\begin{proof}[Derivation of \cref{eq:wealth-jensen-info}]
  We have 
  \begin{align*}
    I_{\phi_{\mathrm{eq}}}(\pop) &= \sum_{i = 1}^n \mu_i \paren{\norm{\mC(\vy_i - \vyref)}^2 - \norm{\mC(\bar{\vy} - \vyref)}^2} \\
    &= \sum_{i = 1}^n \mu_i \paren{\norm{\mC\vy_i}^2 - 2\inner{\mC\vy_i}{\mC\vyref} + \norm{\mC\vyref}^2} \\ 
    &\quad - \sum_{i = 1}^n \mu_i\paren{\norm{\mC\bar{\vy}}^2 + 2\inner{\mC\bar{\vy}}{\mC\vyref} - \norm{\mC\vyref}^2} \\
    &= \sum_{i = 1}^n \mu_i \paren{\norm{\mC\vy_i}^2 - 2\inner{\mC\vy_i}{\mC\vyref} - \norm{\mC\bar{\vy}}^2 + 2\inner{\mC\bar{\vy}}{\mC\vyref}} \\
    &= \sum_{i = 1}^n \mu_i \paren{\norm{\mC\vy_i}^2 - 2\inner{\mC(\vy_i - \bar{\vy})}{\mC\vyref} - \norm{\mC\bar{\vy}}^2 }\\
    &= \paren{\sum_{i = 1}^n \mu_i \norm{\mC\vy_i}^2} - \norm{\mC\bar{\vy}}^2 \;,
  \end{align*}
  since $\sum_{i = 1}^n \mu_i \vy_i = \bar{\vy}$. 
\end{proof}